\chapter{Myocardial Infarction}

\section*{Definition and Overview}
Myocardial infarction (MI, heart attack) is the irreversible necrosis of heart muscle secondary to prolonged ischaemia, most commonly due to acute thrombotic occlusion of a coronary artery following atherosclerotic plaque rupture or erosion. It is classified by aetiology (Type 1 spontaneous, Type 2 supply-demand mismatch, Type 3 sudden cardiac death, Type 4/5 procedural) and by ECG pattern (ST-elevation MI [STEMI] vs non-ST-elevation MI [NSTEMI]). STEMI implies complete occlusion requiring immediate reperfusion; NSTEMI implies partial or transient occlusion. MI is a leading cause of death and disability globally, but outcomes have improved dramatically with timely reperfusion, evidence-based medical therapy, and secondary prevention.

\section*{Epidemiology}
In Australia, approximately 55,000 acute coronary syndrome (ACS) events occur annually, with ~25,000 MIs. Age-standardised rates have declined by >50\% since the 1980s due to risk factor control and treatment advances, but the absolute number is rising with population ageing. Men are affected ~2x more than women before age 70; rates equalise thereafter. Indigenous Australians have 2--3x higher rates and earlier onset (median age 55 vs 70). One-year mortality post-MI is ~10--15\%, highest in the first 30 days. Cardiovascular disease remains the leading cause of death in Australia (~43,000 deaths/year).

\section*{Aetiology and Pathophysiology}
\begin{itemize}
    \item \textbf{Atherosclerosis:} lipid accumulation, inflammation, smooth muscle proliferation, fibrous cap formation; vulnerable plaques have large lipid core, thin fibrous cap, macrophage infiltration
    \item \textbf{Plaque rupture/erosion:} exposes thrombogenic collagen and tissue factor $\to$ platelet adhesion/activation $\to$ thrombus formation $\to$ coronary occlusion
    \item \textbf{Type 1 MI:} spontaneous plaque rupture/erosion in context of atherosclerosis
    \item \textbf{Type 2 MI:} supply-demand mismatch (tachycardia, hypotension, anaemia, hypoxia, severe hypertension, coronary spasm, embolism) without primary plaque event
    \item \textbf{Ischaemic cascade:} within seconds: diastolic dysfunction; minutes: systolic dysfunction, ECG changes; 20--40 min: irreversible necrosis (wavefront from subendocardium to epicardium); hours: completed infarction, scar formation over weeks
    \item \textbf{Reperfusion injury:} oxidative stress, calcium overload, inflammation, mitochondrial permeability transition on restoring flow
\end{itemize}

\section*{Clinical Features}
\begin{itemize}
    \item \textbf{Typical:} central/retrosternal chest pressure/heaviness/squeezing $\ge$20 min, radiating to jaw, neck, left arm, back; associated with diaphoresis, nausea, dyspnoea, sense of impending doom
    \item \textbf{Atypical:} epigastric pain, isolated dyspnoea, fatigue, syncope, palpitations, cardiac arrest; more common in women, elderly, diabetics
    \item \textbf{Silent MI:} no symptoms; detected incidentally on ECG/imaging; more common in diabetes, elderly
    \item \textbf{Physical exam:} may be normal; tachycardia/bradycardia, hypotension/hypertension, cool peripheries, raised JVP, basal crackles (LV failure), new murmur (mitral regurgitation, VSR), pericardial friction rub (Dressler syndrome)
    \item \textbf{Complications presenting acutely:} cardiogenic shock, acute pulmonary oedema, arrhythmias (VF, VT, complete heart block), mechanical complications (papillary muscle rupture, VSR, free wall rupture)
\end{itemize}

\section*{Differential Diagnosis}
\begin{itemize}
    \item \textbf{Unstable angina:} ischaemic symptoms at rest/minimal exertion, no troponin elevation
    \item \textbf{Non-cardiac chest pain:} GERD, musculoskeletal, costochondritis, herpes zoster, anxiety/panic
    \item \textbf{Pulmonary embolism:} pleuritic pain, dyspnoea, tachycardia, hypoxia, RV strain on ECG/echo
    \item \textbf{Aortic dissection:} tearing interscapular pain, pulse/BP differential, widened mediastinum
    \item \textbf{Pericarditis:} positional pleuritic pain, diffuse ST elevation, PR depression, friction rub
    \item \textbf{Myocarditis:} viral prodrome, troponin elevation, normal coronaries, LGE on CMR
    \item \textbf{Takotsubo cardiomyopathy:} apical ballooning post-stress, troponin modestly elevated, normal coronaries
\end{itemize}

\section*{When to Seek Medical Care}
Any chest pain or equivalent symptom $\ge$10--15 min at rest, or new/worsening angina pattern, warrants immediate emergency assessment. Do not drive; call ambulance.

\textbf{Call 000 for:}
\begin{itemize}
    \item Chest pain/discomfort $\ge$10 min, severe, or with diaphoresis, dyspnoea, nausea, radiation
    \item Syncope, presyncope, palpitations with chest pain
    \item Known coronary disease with change in angina pattern (more frequent, less exertion, at rest)
    \item Post-MI: new/worsening chest pain, dyspnoea, swelling, syncope
\end{itemize}

\section*{Diagnosis and Investigations}
\begin{itemize}
    \item \textbf{ECG (within 10 min of presentation):}
    \begin{itemize}
        \item STEMI: ST elevation $\ge$1 mm (limb) or $\ge$2 mm (precordial) in $\ge$2 contiguous leads; new LBBB; posterior MI (ST depression V1--V3 + tall R)
        \item NSTEMI: ST depression, T inversion, or normal/non-specific
        \item Serial ECGs essential
    \end{itemize}
    \item \textbf{Troponin (hs-cTnI/T):} gold standard biomarker; detectable within 1--3 h, peaks 12--24 h; rising/falling pattern confirms acute MI; interpret with clinical context (Type 2 MI, myocarditis, PE, CKD can elevate)
    \item \textbf{0/1-h or 0/2-h accelerated diagnostic protocols:} rule-out/rule-in pathways using hs-troponin + clinical risk scores (HEART, EDACS)
    \item \textbf{Echocardiography:} regional wall motion abnormalities (immediate), LV function, complications (thrombus, VSR, tamponade)
    \item \textbf{Coronary angiography:} definitive anatomy; STEMI $\to$ immediate primary PCI; NSTEMI $\to$ early invasive (within 24 h high-risk, 72 h intermediate)
    \item \textbf{Other:} FBC, lipids, HbA1c, renal function, electrolytes, coagulation, CK-MB (if troponin unavailable)
\end{itemize}

\section*{Management}
\begin{itemize}
    \item \textbf{Pre-hospital / immediate (STEMI):}
    \begin{itemize}
        \item Aspirin 300 mg chewed, P2Y12 inhibitor (ticagrelor 180 mg or clopidogrel 600 mg), anticoagulant (unfractionated heparin or enoxaparin or bivalirudin)
        \item Primary PCI within 90 min (door-to-balloon) $\to$ gold standard; if PCI not available within 120 min $\to$ fibrinolysis (tenecteplase) then transfer for PCI
        \item Oxygen if SpO2 $<$90\%; nitrate for pain/BP; morphine if needed; avoid routine oxygen
    \end{itemize}
    \item \textbf{NSTEMI:} early invasive strategy (angiography within 24 h high-risk [GRACE >140, dynamic ECG, rising troponin, HF, arrhythmia]; 72 h others); medical stabilisation first
    \item \textbf{Medical therapy (all MI, start early, continue lifelong):}
    \begin{itemize}
        \item Dual antiplatelet therapy (DAPT): aspirin 100 mg daily + ticagrelor 90 mg BD (preferred) or clopidogrel 75 mg daily $\times$ 12 months (shorter if bleeding risk, longer if high ischaemic risk)
        \item High-intensity statin (atorvastatin 40--80 mg or rosuvastatin 20--40 mg) -- LDL target $<$1.4 mmol/L or $>$50\% reduction
        \item ACE inhibitor (or ARB) $\to$ start early, titrate to max tolerated; ARNI (sacubitril/valsartan) if HFrEF
        \item Beta-blocker $\to$ start early if no contraindication; continue indefinitely post-MI with reduced EF
        \item Mineralocorticoid receptor antagonist (eplerenone) if EF $\le$40\% + HF/diabetes
        \item SGLT2 inhibitor (dapagliflozin/empagliflozin) -- Class I for HFrEF, now recommended post-MI regardless of EF/diabetes
        \item Consider icosapent ethyl if TG elevated on statin
    \end{itemize}
    \item \textbf{Mechanical complications:} urgent surgery (VSR, free wall rupture, papillary muscle rupture)
    \item \textbf{Cardiac rehabilitation:} Phase II (supervised exercise, education, risk factor modification) -- reduces mortality, readmission; refer all
    \item \textbf{Secondary prevention:} smoking cessation, Mediterranean diet, weight management, physical activity, BP $<$130/80, diabetes control, influenza/pneumococcal vaccination, mental health screening
\end{itemize}

\section*{Complications}
\begin{itemize}
    \item \textbf{Early:} arrhythmias (VF/VT, bradyarrhythmias), cardiogenic shock, acute HF/pulmonary oedema, mechanical (VSR, free wall rupture, papillary muscle rupture), pericarditis, thromboembolism
    \item \textbf{Late:} chronic HF (ischaemic cardiomyopathy), ventricular aneurysm, recurrent MI, angina, depression/anxiety, post-MI pericarditis (Dressler)
\end{itemize}

\section*{Prognosis and Outcomes}
Short-term mortality has fallen to ~5--7\% (in-hospital) with primary PCI and modern therapy. Long-term prognosis depends on infarct size, LV function (EF), completeness of revascularisation, adherence to secondary prevention, and comorbidities. Five-year survival ~75--80\% for those discharged alive. Cardiac rehabilitation participation reduces all-cause mortality by 20--30\%. Quality of life is generally good with optimal management.

\section*{Prevention and Self-Care}
\begin{itemize}
    \item \textbf{Primary prevention:} assess absolute CVD risk (Australian calculator) from age 45 (35 Indigenous); treat risk factors
    \item \textbf{Lifestyle:} Mediterranean diet; 150--300 min/week moderate activity; healthy weight; no smoking; limit alcohol
    \item \textbf{Risk factor targets:} BP $<$130/80; LDL $<$1.8 mmol/L (very high risk) or $<$1.4 mmol/L (extreme risk); HbA1c $\le$7\%; non-HDL $<$2.2 mmol/L
    \item \textbf{Medication adherence:} DAPT, statin, ACEi/ARB, beta-blocker, MRA, SGLT2i -- lifelong unless contraindicated
    \item \textbf{Cardiac rehab:} attend and complete program
    \item \textbf{Know warning signs:} chest pain action plan; keep GTN spray accessible
    \item \textbf{Regular review:} GP 6-weekly initially, then 3--6 monthly; cardiologist annually
    \item \textbf{Vaccinations:} annual influenza, pneumococcal, COVID boosters
\end{itemize}

\section*{Sources and Further Reading}
The following reputable sources provide further information for healthcare students and patients seeking reliable, current advice about myocardial infarction.

\begin{itemize}
    \item National Heart Foundation of Australia / Cardiac Society of Australia and New Zealand. \textit{Acute coronary syndromes guidelines.} https://www.heartfoundation.org.au/
    \item Australian and New Zealand Committee on Resuscitation (ANZCOR). \textit{ACS guidelines.} https://resus.org.au/
    \item eviQ Cancer Treatments Online. \textit{Cardiotoxicity protocols (relevant for cancer patients).} https://www.eviq.org.au/
    \item Healthdirect Australia. \textit{Heart attack.} https://www.healthdirect.gov.au/heart-attack
    \item Better Health Channel. \textit{Heart attack.} https://www.betterhealth.vic.gov.au/
    \item European Society of Cardiology. \textit{STEMI and NSTEMI guidelines.} https://www.escardio.org/
    \item American Heart Association / American College of Cardiology. \textit{MI guidelines.} https://www.heart.org/
    \item Australian Cardiovascular Health and Rehabilitation Association (ACRA). \textit{Cardiac rehabilitation standards.} https://acra.org.au/
\end{itemize}